\section{Реализация}

В данном разделе описывается программная реализация исследования на языке \textit{Python}.
Весь код структурирован в виде Jupyter Notebook и разбит на логические блоки:
предобработка данных, реализация методов отбора, оценка параметров выборки и симуляция.

% ─────────────────────────────────────────────────────────────────────────────
\subsection{Используемые библиотеки}

Для реализации алгоритмов использовались следующие библиотеки:

\begin{itemize}[noitemsep]
    \item \texttt{numpy} --- векторизованные вычисления и генерация псевдослучайных чисел;
    \item \texttt{pandas} --- загрузка, хранение и агрегация табличных данных;
    \item \texttt{scipy.stats} --- квантили распределений Стьюдента и нормального;
    \item \texttt{matplotlib} и \texttt{seaborn} --- визуализация результатов;
    \item \texttt{functools.partial} --- частичное применение функций для единообразного вызова методов отбора.
\end{itemize}

% ─────────────────────────────────────────────────────────────────────────────
\subsection{Загрузка и предобработка данных}

Исходные данные загружаются из файла Excel и сортируются по коду респондента. Для каждого респондента формируется метка страты как конкатенация его возрастной группы и пола:

\begin{mintedbox}{python}
df["Stratum"] = df["Возрастная група"] + "_" + df["Пол"]
\end{mintedbox}

Столбцы мастер-таблицы разбиваются на блоки: метаданные, оценки \textit{идеального} продукта, веса \textit{важности} параметров и оценки каждого из трeх брендов (Sony, JBL, Marshall).

\begin{mintedbox}{python}
data_cols = data.columns[METADATA_COLS:]
chunks = np.split(data_cols, SECTIONS_COUNT)
ideal_cols, imp_cols, *brand_col_groups = chunks
\end{mintedbox}

% ─────────────────────────────────────────────────────────────────────────────
\subsection{Вычисление интегральных показателей по модели идеальной точки}

Для каждого респондента и каждого бренда вычисляется интегральный балл отклонения от идеала по формуле~\ref{eq:integral_score} из раздела методологии.

Реализация использует векторизованные операции \texttt{NumPy}:

\begin{mintedbox}{python}
for name in BRAND_NAMES:
    brand_cols = brands[name]
    score = (np.abs(ideal_col - df[brand_cols]) * imp_col).sum(axis=1)
    all_scores.append(score)

pop_scores = np.column_stack(all_scores)
\end{mintedbox}

Результатом является матрица \texttt{scores} размером $N \times 3$, где каждая строка соответствует респонденту, а столбцы --- трeм брендам. На еe основе вычисляются \textit{истинные} (генеральные) средние значения и определяется бренд, наиболее близкий к идеалу.

\medskip
\noindent Результаты по генеральной совокупности:

\begin{table}[h!]
\centering
\begin{tblr}{
    colspec = {|X[l]|X[c]|},
    row{1} = {gray!30, font=\bfseries},
    row{even} = {gray!10},
    hlines,
    vlines,
}
Бренд & Среднее отклонение от идеальной точки \\
Sony     & 76{,}2565 \\
JBL      & 76{,}3946 \\
Marshall & 76{,}9517 \\
\end{tblr}
\caption{Генеральные средние по модели идеальной точки}
\label{tab:general_means}
\end{table}

\noindent Бренд, наиболее близкий к идеалу --- \textbf{Sony}.

% ─────────────────────────────────────────────────────────────────────────────
\subsection{Реализация методов отбора}

Для уменьшения времени выполнения кода в дальнейшем было решено реализовать функции отбора таким образом, чтобы они возвращали индексы строк из таблицы, а не сами строки.

Несмотря на то что \textit{механический} отбор подразумевает только бесповторный характер, для унификации все функции отбора обладают тремя обязательными общими параметрами: \texttt{n} --- желаемым объемом выборки, \texttt{replace}~--- характером отбора (повторный или бесповторный) и \texttt{seed} --- зерном для воспроизводимости результатов (в случае \textit{механического} отбора этот параметр назван \texttt{start\_iter} --- начальная точка отсчета).

\subsubsection*{Случайный отбор}

Выборка формируется с помощью генератора псевдослучайных чисел\\
\texttt{numpy.random.default\_rng} путeм случайного выбора $n$ индексов из генеральной совокупности (с повторением или без):

\begin{mintedbox}{python}
def sample_random(n_total: int, n: int, replace: bool, seed: int) -> np.ndarray:
    indices_random = np.random.default_rng(seed).choice(n_total, size=n, replace=replace)
    return indices_random
\end{mintedbox}

\subsubsection*{Механический отбор}

Элементы отбираются через равный шаг $k = \lfloor N / n \rfloor$, начиная с фиксированной стартовой позиции в середине массива. Данный метод поддерживает только бесповторный способ (несмотря на параметр \texttt{replace} --- он ни на что не влияет):

\begin{mintedbox}{python}
def sample_mechanical(n_total: int, n: int, replace: bool = False, start_iter: int = 0) -> np.ndarray:
    k = N_TOTAL // n
    start = N_TOTAL // 2 + start_iter
    indices_mech = (start + np.arange(n) * k) % n_total
    return indices_mech
\end{mintedbox}

\subsubsection*{Типический (стратифицированный) отбор}

Генеральная совокупность предварительно разбита на страты по \\
поло-возрастному признаку. Из каждой страты случайным образом отбирается количество
элементов, пропорциональное еe доле в совокупности:

\clearpage

\begin{mintedbox}{python}
def sample_typical(stratum_weights: dict, n: int, replace: bool, seed: int, get_strats: bool = False):
    rng = np.random.default_rng(seed)
    parts = []
    
    for s, w in stratum_weights.items():
        pool = strata_groups[s]
        n_s = max(1, round(n * w))
        n_s = n_s if replace else min(n_s, len(pool))
        parts.append(rng.choice(pool, size=n_s, replace=replace))

    indices_typ = np.concatenate(parts)

    if not get_strats:
        return indices_typ
    
    return indices_typ, parts
\end{mintedbox}

\subsubsection*{Кластерный отбор}

Единицей отбора выступает кластер (станция). Количество отбираемых кластеров вычисляется исходя из средней наполненности кластера. Затем случайным образом выбираются $r$ кластеров, и все входящие в них респонденты включаются в выборку:

\begin{mintedbox}{python}
def sample_cluster(cluster_keys: list, n: int, replace: bool, seed: int, get_clusters: bool = False):
    rng = np.random.default_rng(seed)
    r = round(n / AVG_PEOPLE_PER_CLUSTER)
    r = r if replace else min(r, len(cluster_keys))
    
    chosen_clusters = rng.choice(cluster_keys, size=r, replace=replace)
    selected = [cluster_groups[c] for c in chosen_clusters]
    indices_cluster = np.concatenate(selected)

    if not get_clusters:
        return indices_cluster

    return indices_cluster, selected
\end{mintedbox}

% ─────────────────────────────────────────────────────────────────────────────
\subsection{Оценка дисперсий разведочной выборки}

На основе разведочной выборки объeма\footnote{Согласно техническому заданию объем разведочной совокупности принимается равным $4\%$ от генеральной совокупности.} оцениваются дисперсии, соответствующие каждому методу отбора):

\begin{itemize}
    \item \textbf{Случайный и механический} --- общая выборочная дисперсия~\eqref{eq:var},
    \item \textbf{Типический} --- средняя из внутригрупповых дисперсий страт~\eqref{eq:var_within_mean},
    \item \textbf{Кластерный} --- межгрупповая дисперсия средних по кластерам~\eqref{eq:var_between}.
\end{itemize}

\begin{mintedbox}{python}
leader_scores = pop_scores[:, best_brand]

# Случайный
idx_rnd = methods["Случайный"](N_SAMPLE, True, SEED)
var_rnd = leader_scores[idx_rnd].var(ddof=1)

# Механический
idx_mech = methods["Механический"](N_SAMPLE, False, SEED)
var_mech = leader_scores[idx_mech].var(ddof=1)

# Типический
idx_typ, parts = methods["Типический"](N_SAMPLE, True, SEED, get_strats=True)
sizes = np.array([len(p) for p in parts])
variances = np.array([leader_scores[p].var(ddof=1) for p in parts])
var_typ = np.average(variances, weights=sizes)

# Кластерный
idx_clst, selected = methods["Кластерный"](N_SAMPLE, True, SEED, get_clusters=True)
cluster_means = np.array([leader_scores[p].mean() for p in selected])
var_clst = np.var(cluster_means)
\end{mintedbox}
